
\subsubsection*{5.1.1}
\phantomsection
\label{sga3:VIIB:5.1.1}

Consider the following category \(\mathcal C\). An object of
\(\mathcal C\) is a pair \((A,J)\) consisting of a profinite
\(k\)-algebra \(A\) and a closed ideal \(J\subset A\). A morphism
\[
  \psi:(A,J)\longrightarrow(A',J')
\]
in \(\mathcal C\) is a continuous \(k\)-algebra homomorphism
\[
  A\longrightarrow A'
\]
that maps \(J\) into \(J'\).

The assignment
\[
  (A,J)
  \longmapsto
  \bigl(\operatorname{Spf}(A/J),\operatorname{Spf}(A)\bigr)
\]
is an antiequivalence from \(\mathcal C\) to the category of pairs
\((Z,Y)\) consisting of a formal \(k\)-variety \(Y\) and a formal
subvariety \(Z\subset Y\). A morphism
\[
  \varphi:(Z,Y)\longrightarrow(Z',Y')
\]
is a morphism of formal \(k\)-varieties
\[
  Y\longrightarrow Y'
\]
that maps \(Z\) into \(Z'\).

A cogroup structure on an object \((A,J)\) of \(\mathcal C\)
consists of a formal-group structure on
\(\operatorname{Spf}(A)\) satisfying the following conditions
\textup{(with the notation of \SGARefLink{sga3:VIIB:2.1}{2.1})}:
\[
\begin{aligned}
  \textup{(1)}\quad&
  \Delta_A(J)
  \subset
  J\widehat\otimes_kA
  +
  A\widehat\otimes_kJ,
  \\
  \textup{(2)}\quad&
  \varepsilon_A(J)=0,
  \\
  \textup{(3)}\quad&
  c_A(J)\subset J.
\end{aligned}
\]
These conditions also say that
\[
  H=\operatorname{Spf}(A/J)
\]
is a formal subgroup of
\[
  G=\operatorname{Spf}(A).
\]
\footnote{\label{sga3:VIIB:nde:145}If \((A',J')\) is a second
cogroup of \(\mathcal C\), corresponding to a pair
\(H'\subset G'\) of formal \(k\)-groups, then giving a cogroup
morphism
\[
  (A',J')\longrightarrow(A,J)
\]
is equivalent to giving a morphism of formal \(k\)-groups
\[
  G\longrightarrow G'
\]
that maps \(H\) into \(H'\).}

Suppose in addition that \(A\) is local; equivalently,
\(\operatorname{Spf}(A)\) is an infinitesimal formal group. If
\(J\ne A\), then conditions \textup{(2)} and \textup{(3)} follow
from \textup{(1)}. Indeed, a proper closed ideal \(J\) of \(A\)
is contained in the augmentation ideal \(I_A\), so
\textup{(2)} holds. Moreover,
\[
  M=\operatorname{Spf}(A/J)
\]
is a formal submonoid of \(G\). Since \(G\) is infinitesimal,
\SGARefLink{sga3:VIIB:2.7:prop}{Proposition~2.7} shows that \(M\) is a formal subgroup of \(G\);
hence \textup{(3)} also holds.
